\chapter{小数计算}

\zsd{知识要点}


 (1) 计算整数时所学的所有方法都可以利用到小数中来.

 (2) 注意小数点的位置.

 (3) 小数点的移动: 左移一位, 缩小 10 倍, 右移一位, 扩大 10 倍.

 
 \zsd{典例精讲}
  
 
\subzsd{加减运算}
 %%%例题%%%%%%%%%%%%%%%
\begin{example}
	\begin{problems}
		计算(1)$7.973+1.275-1.473+2.225$ \hfil   (2)$ 4.75-9.64+(8.25-1.36)$ 
	\end{problems} 
	
	\begin{Ssol}
		\vspace{\solfourproblems}
	\end{Ssol}
\end{example}

 %%%练习%%%%%%%%%%%%%%%
\begin{lianxi}
	\begin{problems}
		计算：（1） $1.34+34.1+2.56$ ；（2） $2.1+2.3+2.5+2.7+2.9$ ．
		
	\end{problems} 
	
	\begin{lxSsol}
		\vspace{\solthreeproblems}
	\end{lxSsol}
\end{lianxi}
 %%%例题%%%%%%%%%%%%%%%
\begin{example}
\begin{problems}
		计算: $12.34+23.45+34.56+45.67+56.78+67.89+78.91+89.12+91.23$
\end{problems} 

	
	\begin{Ssol}
		\vspace{\solfourproblems}
	\end{Ssol}
\end{example}




\subzsd{乘除法巧算}

 
%%%例题%%%%%%%%%%%%%%%
\begin{example}
	\begin{problems}
		计算：（1） $1.23 \times 100$ ；（2） $3.56 \times 4$ ；（3） $2.5 \times 0.4$ ；（4） $1.23 \div 100$ ；（5） $0.84 \div 12$ ；（6） $1.771 \div 1.1$ ．
		
	\end{problems} 
	
	\begin{Ssol}
		\vspace{\solthreeproblems}
	\end{Ssol}
\end{example}


%%%练习%%%%%%%%%%%%%%%
\begin{lianxi}
	\begin{problems}
		计算：（1） $123.45 \times 0.01$ ；（2） $44.5 \times 0.6$ ；（3） $123.45 \div 0.01$ ；（4） $1 \div 8$ ．
	\end{problems} 
	
	\begin{lxMsol}
		\vspace{\solthreeproblems}
	\end{lxMsol}
\end{lianxi}


%%%例题%%%%%%%%%%%%%%%
\begin{example}
	\begin{problems}
		计算：（1） $1.25 \times 0.26 \times 8$ ；（2） $4.5 \times 4.8 \div 15 \div 0.24$ ．
		
	\end{problems} 
	
	\begin{Ssol}
		\vspace{\solthreeproblems}
	\end{Ssol}
\end{example}


%%%练习%%%%%%%%%%%%%%%
\begin{lianxi}
	\begin{problems}
		计算：（1） $0.25 \times 1.36 \times 40$ ；（2） $0.56 \times 2.3 \div 1.4$ ．
		
	\end{problems} 
	
	\begin{lxMsol}
		\vspace{\solthreeproblems}
	\end{lxMsol}
\end{lianxi}
%%%%%%%%%%%%%%%%%%
 %%%例题%%%%%%%%%%%%%%%
\begin{example}
	计算: $0.125 \times 0.25 \times 0.5 \times 64$

	
	\begin{Ssol}
		\vspace{\solfourproblems}
	\end{Ssol}
\end{example}


%%%%%%%%%%%%%%%%%%

 %%%例题%%%%%%%%%%%%%%%
\begin{example}
计算 $(3.9 \times 5.5 \times 6.3 \times 3.6) \div(1 \times 3 \times 5 \times 7 \times 0.9 \times 1.1 \times 1.3)$	
	
	\begin{Ssol}
		\vspace{\solfourproblems}
	\end{Ssol}
\end{example}

 %%%例题%%%%%%%%%%%%%%%
\begin{example}
	\begin{problems}
		已知 $1.08 \div 1.2 \div 2.3=10.8 \div \square$ ，其中 $\square$ 表示的数是\tk.
	\end{problems} 
	
	\begin{Ssol}
		\vspace{\solfourproblems}
	\end{Ssol}
\end{example}




\subzsd{运算律}

 %%%例题%%%%%%%%%%%%%%%
\begin{example}
	\begin{problems}
		计算: $34.5 \times 8.23-34.5+2.77 \times 34.5$
	\end{problems} 
	
	\begin{Ssol}
		\vspace{\solfourproblems}
	\end{Ssol}
\end{example}


%%%练习%%%%%%%%%%%%%%%
\begin{lianxi}
	\begin{problems}
		计算： $4.5 \times 7-4.5 \times 6.94$ ．
	\end{problems} 
	
	\begin{lxMsol}
		\vspace{\solthreeproblems}
	\end{lxMsol}
\end{lianxi}



 %%%例题%%%%%%%%%%%%%%%
\begin{example}
计算: (1) $8.6 \times 0.37+0.73 \times 8.6$ ．

 (2)$6.25 \times 0.16+264 \times 0.0625+5.2 \times 6.25+0.625 \times 20$    
	\begin{Ssol}
		\vspace{\solfourproblems}
	\end{Ssol}
\end{example}
 %%%例题%%%%%%%%%%%%%%%

 %%%练习%%%%%%%%%%%%%%%
\begin{lianxi}
	\begin{problems}
		 $4.32 \times 23.5+0.581 \times 568+43.2 \times 3.46$ ．
	\end{problems} 
	
	\begin{lxMsol}
		\vspace{\solthreeproblems}
	\end{lxMsol}
\end{lianxi}


 %%%例题%%%%%%%%%%%%%%%
\begin{example}
计算：
	$1.24 \times 20+124 \times 0.7+0.248 \times 550$
	
	
	\begin{Ssol}
		\vspace{\solfourproblems}
	\end{Ssol}
\end{example}


%%%%%%%%%%%%%%%%%%
%%%%%%%%%%%%%%%%%%
 %%%例题%%%%%%%%%%%%%%%
\begin{example}
计算: $1.25 \times 67.875+25 \times 33.9375+1250 \times 0.053375$
	
	
	
	\begin{Ssol}
		\vspace{\solfourproblems}
	\end{Ssol}
\end{example}

 %%%例题%%%%%%%%%%%%%%%
\begin{example}
	\begin{problems}
		计算: $13.5 \times 9.9+6.5 \times 10.1$
	\end{problems} 
	
	\begin{Ssol}
		\vspace{\solthreeproblems}
	\end{Ssol}
\end{example}


%%%%%%%%%%%%%%%%%%


\subzsd{提高例题}

 %%%例题%%%%%%%%%%%%%%%
\begin{example}
	计算： $0.2+0.5+0.8+0.11+0.14+\cdots+0.101$	
	
	\begin{Ssol}
		\vspace{\solfourproblems}
	\end{Ssol}
\end{example}


 %%%例题%%%%%%%%%%%%%%%
\begin{example}
计算：	$1 \times 1.9+2 \times 1.8+\cdots+18\times 0.2+19 \times 0.1$
	
	\begin{Ssol}
		\vspace{\solfourproblems}
	\end{Ssol}
\end{example}


%
%%%%%%%例题%%%%%%%%%%%%%%%%%%%%%%%%%%%%%%
%\begin{example}
%从 $A$ 点走到 $D$ 点, 要求必须经过 $B$ 点, 不能经过 $C$ 点, 只能向右和向下走, 那 么有多少种走法?
%	\begin{figure}[H]
%		\begin{flushright}
%			\begin{tikzpicture}[scale=0.6]
%			\draw (0,0) grid (7,5);
%			\draw[fill=blue] (0,5) circle(1.5pt);
%			\draw[fill=blue] (2,3) circle(1.5pt);
%			\draw[fill=blue] (4,2) circle(1.5pt);
%			\draw[fill=blue] (7,0) circle(1.5pt);
%			
%			\node[left]  at (0,5) {$A$};
%			\node[left]  at (2,3) {$B$};
%			\node[left]  at (4,2) {$C$};
%			\node[below]  at (7,0) {$D$};
%			\end{tikzpicture}
%		\end{flushright}
%		
%	\end{figure}
%	\begin{sol}
%		对于计算行走路线的问题,“标数法”是一种非常有效的方法．标数过程中要注意“必过点”和“不过点”．
%		\vspace{\solfiveproblems}
%		
%	\end{sol}
%\end{example}




